\chapter*{Úvod}
\phantomsection
\addcontentsline{toc}{chapter}{Úvod}
Projekt Deltabot vznikl v rámci projektu v soutěží FEKTTeams 2024. Na různých částech
projektu pracoval čtyřčlenný tým, jejichž práce je popsána v jejich bakalářských pracích. 

Tato práce navazuje na bakalářskou práci Řízení pohybu DeltaBotu, která se zabývala
zprovozněním robotického systému vlastní konstrukce pomocí PLC od firmy B\&R.
Zároveň využívá analytické modely výpočtů polohy jednotlivých pozic motorů na ose
a koncového bodu robotu. Těmito modely se zabývala práce Analytické modelování delta
robota.

% Kdyžtak dopsat
Práce se zabývá problematikou kolizí robotů, jejichž pracovní oblasti se protínají.
Výsledkem práce by mělo být jak kolize zjistit a vyvarovat se jim. Existují tři hlavní
myšlenky, jak by se mohly kolize robotů řešit.

V rámci práce byl stávající PLC projekt převeden na novější verzi, kde je od základu
nastavena bezpečná komunikace na síti.

 
% no toto tam možná ani nebude protože to vlastne ke své diplomce nepotřebuji XD

Jelikož pro řešení kolizí je potřeba dva roboty, které fyzicky nebyly k dispozici,
vše je řešeno přes simulaci. Podstatnou částí semestrální práce proto je vytvoření
simulovaného prostředí se dvěma roboty. Tato část následně může být rozšířena o
simulaci řešení rovnoměrného zatížení těchto robotů, což by mělo být splněno v rámci
celé diplomové práce. 




%Úvod studentské práce, např\,\dots

%Nečíslovaná kapitola Úvod obsahuje \uv{seznámení} čtenáře s~problematikou práce.
%Typicky se zde uvádí:
%(a) do jaké tematické oblasti práce spadá, (b) co jsou hlavní cíle celé práce a (c) jakým způsobem jich bylo dosaženo.
%Úvod zpravidla nepřesahuje jednu stranu.
%Poslední odstavec Úvodu standardně představuje základní strukturu celého dokumentu.

%Tato práce se věnuje oblasti \acs{DSP} (\acl{DSP}), zejména jevům, které nastanou při nedodržení Nyquistovy podmínky pro \ac{symfvz}.%
%\footnote{Tato věta je pouze ukázkou použití příkazů pro sazbu zkratek.}

%Šablona je nastavena na \emph{dvoustranný tisk}.
%Nebuďte překvapeni, že ve vzniklém PDF jsou volné stránky.
%Je to proto, aby důležité stránky jako např.\ začátky kapitol začínaly po vytisknutí a svázání vždy na pravé straně.
%
%Pokud máte nějaký závažný důvod sázet (a~zejména tisknout) jednostranně, nezapomeňte si přepnout volbu \texttt{twoside} na \texttt{oneside}!